% Options for packages loaded elsewhere
\PassOptionsToPackage{unicode}{hyperref}
\PassOptionsToPackage{hyphens}{url}
\PassOptionsToPackage{dvipsnames,svgnames,x11names}{xcolor}
%
\documentclass[
  10pt,
  a4paper,
]{article}
\usepackage{amsmath,amssymb}
\usepackage{setspace}
\usepackage{iftex}
\ifPDFTeX
  \usepackage[T1]{fontenc}
  \usepackage[utf8]{inputenc}
  \usepackage{textcomp} % provide euro and other symbols
\else % if luatex or xetex
  \usepackage{unicode-math} % this also loads fontspec
  \defaultfontfeatures{Scale=MatchLowercase}
  \defaultfontfeatures[\rmfamily]{Ligatures=TeX,Scale=1}
\fi
\usepackage{lmodern}
\ifPDFTeX\else
  % xetex/luatex font selection
  \setmainfont[]{Liberation Serif}
  \setsansfont[]{Liberation Sans}
  \setmonofont[]{DejaVu Sans Mono}
\fi
% Use upquote if available, for straight quotes in verbatim environments
\IfFileExists{upquote.sty}{\usepackage{upquote}}{}
\IfFileExists{microtype.sty}{% use microtype if available
  \usepackage[]{microtype}
  \UseMicrotypeSet[protrusion]{basicmath} % disable protrusion for tt fonts
}{}
\makeatletter
\@ifundefined{KOMAClassName}{% if non-KOMA class
  \IfFileExists{parskip.sty}{%
    \usepackage{parskip}
  }{% else
    \setlength{\parindent}{0pt}
    \setlength{\parskip}{6pt plus 2pt minus 1pt}}
}{% if KOMA class
  \KOMAoptions{parskip=half}}
\makeatother
\usepackage{xcolor}
\usepackage[margin=22mm]{geometry}
\setlength{\emergencystretch}{3em} % prevent overfull lines
\providecommand{\tightlist}{%
  \setlength{\itemsep}{0pt}\setlength{\parskip}{0pt}}
\setcounter{secnumdepth}{-\maxdimen} % remove section numbering
\ifLuaTeX
\usepackage[bidi=basic]{babel}
\else
\usepackage[bidi=default]{babel}
\fi
\babelprovide[main,import]{korean}
\ifPDFTeX
\else
\babelfont{rm}[]{Liberation Serif}
\fi
% get rid of language-specific shorthands (see #6817):
\let\LanguageShortHands\languageshorthands
\def\languageshorthands#1{}
\usepackage{xeCJK}
\setCJKmainfont{Noto Serif CJK KR}
\setCJKsansfont{Noto Sans CJK KR}
\setCJKmonofont{Noto Sans Mono CJK KR}
\xeCJKsetup{CJKspace=true}
\usepackage{xurl}
\usepackage{fvextra}
\usepackage{fancyhdr}
\AtBeginDocument{\DefineVerbatimEnvironment{verbatim}{Verbatim}{breaklines=true,breakanywhere=true,fontsize=\small}}
\setlength{\emergencystretch}{4em}
\setlength{\headheight}{15pt}
\pagestyle{fancy}\fancyhf{}
\fancyhead[L]{\small 중성자별 구조와 비방사 진동}
\fancyfoot[C]{\thepage}
\renewcommand{\headrulewidth}{0.3pt}
\setcounter{tocdepth}{2}
\ifLuaTeX
  \usepackage{selnolig}  % disable illegal ligatures
\fi
\usepackage{bookmark}
\IfFileExists{xurl.sty}{\usepackage{xurl}}{} % add URL line breaks if available
\urlstyle{same}
\hypersetup{
  pdftitle={A00 · 실행환경과 입력의 의미},
  pdflang={ko-KR},
  colorlinks=true,
  linkcolor={Maroon},
  filecolor={Maroon},
  citecolor={Blue},
  urlcolor={Blue},
  pdfcreator={LaTeX via pandoc}}

\title{A00 · 실행환경과 입력의 의미}
\usepackage{etoolbox}
\makeatletter
\providecommand{\subtitle}[1]{% add subtitle to \maketitle
  \apptocmd{\@title}{\par {\large #1 \par}}{}{}
}
\makeatother
\subtitle{과제 큐시트}
\author{}
\date{}

\begin{document}
\maketitle

\setstretch{1.12}
\textbf{선행:} 없음　 \textbf{독립 작업:} 0.5--1시간 (계획치)

\subsection{프로그램이 끝났다면, 우리는 무엇을 확인한
것일까?}\label{uxd504uxb85cuxadf8uxb7a8uxc774-uxb05duxb0acuxb2e4uxba74-uxc6b0uxb9acuxb294-uxbb34uxc5c7uxc744-uxd655uxc778uxd55c-uxac83uxc77cuxae4c}

아직 별을 계산하지 않았다. 입력한 실수·정수·문자열이 되돌아왔다는 사실로
확인할 수 있는 것 하나와 확인할 수 없는 것 하나를 먼저 적자.

먼저 한두 문장으로 답한 뒤 아래 설명을 읽자. 첫 예측이 틀렸거나 아직
모르겠다고 적는 것은 감점하지 않는다.

\subsection{필요한 관계와 읽을
위치}\label{uxd544uxc694uxd55c-uxad00uxacc4uxc640-uxc77duxc744-uxc704uxce58}

이 과제의 probe는 무차원이다. \(\mathrm{workspace}=\mathrm{learner}\)는
실행 경로의 이름이며 물리 정확도에 대한 진술이 아니다.

\begin{itemize}
\tightlist
\item
  \texttt{START\_HERE.md}
\item
  \texttt{support/fpm.toml}
\item
  \texttt{learner/fpm.toml}
\end{itemize}

\subsection{계산으로
옮기기}\label{uxacc4uxc0b0uxc73cuxb85c-uxc62euxae30uxae30}

계산용 패키지에서 시작해도 된다. 아래 learner 선택은 입력 전달을
확인하는 시험이며 아직 별의 물리를 요구하지 않는다.

\begin{enumerate}
\def\labelenumi{\arabic{enumi}.}
\tightlist
\item
  새 디렉터리에서 입력 파일과 실제 실행할 소스를 확인한다.
\item
  run\_assignment의 --plan으로 시험 경로를 읽고, --run으로 environment
  case를 실행한다.
\item
  입력의 실수·정수·문자열을 바꾸고 실제 전달된 값과 반환 상태를
  비교한다.
\end{enumerate}

\newpage

\subsection{실행해서 무엇을 볼
것인가}\label{uxc2e4uxd589uxd574uxc11c-uxbb34uxc5c7uxc744-uxbcfc-uxac83uxc778uxac00}

패키지 루트에서 실행한다. \texttt{-\/-plan}은 경로만 확인하며 계산
성공을 뜻하지 않는다.

\begin{verbatim}
python tools/run_assignment.py A00 --plan
python tools/run_assignment.py A00 --run \
  --workspace learner --output ../A00_result.json
\end{verbatim}

\subsection{결과를
설명하기}\label{uxacb0uxacfcuxb97c-uxc124uxba85uxd558uxae30}

입력 전달과 실행 대상을 구별했다면 충분하다. 이 시험에서 질량이나 모드의
정확성을 근거로 들었다면, 아직 계산하지 않은 양을 찾아 설명에서 빼자.

\textbf{남길 증거.} 도구 판본과 첫 명령 출력; 빌드 성공이 아직 증명하지
않은 것 두 가지

\textbf{판단 기준.} 2점 예: manifest→source→binary→input의 연결과,
이것이 방정식 검증을 대신하지 않는 이유를 함께 설명한다. 최초 예측의
정오보다, 결과를 본 뒤 이유를 고치고 다른 입력에도 적용하는지를 본다.

\subsection{막혔을 때}\label{uxb9c9uxd614uxc744-uxb54c}

\textbf{힌트 1.} 최초 실패 지점이 Python import, manifest 읽기,
compiler, Fortran 상태 중 어디인지 찾는다.

\textbf{힌트 2.} 상수·별 계산을 호출하기 전에 environment의 입력
문자열과 workspace부터 확인한다.

\subsection{다음 만남에 다시
묻기}\label{uxb2e4uxc74c-uxb9ccuxb0a8uxc5d0-uxb2e4uxc2dc-uxbb3buxae30}

다음 만남에는 로그의 성공 문구를 가리고, 어떤 소스를 어떤 입력으로
실행했는지 파일만 보고 복원한다. 먼저 짧게 답한 뒤 본문과 자신의 기록을
열어 확인한다.

\textbf{범위.} Linux x86\_64에서 확인한 판본만 지원 증거가 있다. 다른
플랫폼 및 실측 학생 효과는 미검증이다.

\end{document}
